\chapter{الأمواج الميكانيكية}\label{ch:g12:mechanical-waves}

مساء السبت، وملعب ممتلئ: تجتاح \omterm{def:g10:signals-and-waves:wave}{موجة} من المتفرجين الواقفين
المدرجاتِ دورةً كاملة في أقل من دقيقة، ومع ذلك لا يغادر أحد
مقعده. وحلقة الحصاة تعبر البركة بينما لا تفعل الفلّينة إلا أن ترتفع وتنخفض. ويدرس
هذا الفصل المسافر نفسه --- أي \omterm{def:g10:signals-and-waves:wave}{الموجة} --- والأعداد
الأربعة التي تروّضها: السرعة \omterm{def:g12:mechanical-waves:delay}{والتأخير} \omterm{def:g10:signals-and-waves:period}{والدور} \omterm{def:g12:mechanical-waves:wavelength}{وطول الموجة} --- وهي تكفي
لقراءة زلزال من شريط ورق.

\section{الاضطراب ينتقل والوسط يبقى}

\begin{definition}[الموجة الميكانيكية]\label{def:g12:mechanical-waves:wave}
\emph{الموجة الميكانيكية}\index{موجة ميكانيكية} اضطراب
ينتشر في وسط مادي --- حبل أو ماء أو هواء أو صخر ---
حاملًا \omterm{def:g10:energy-conservation:energy}{طاقة} \emph{دون نقل مادة}: فكل جزء من
الوسط يتأرجح برهةً في مكانه، ويسلّم الحركة، ثم يستقرّ.
\end{definition}

\begin{example}[لا نقل للمادة]\label{ex:g12:mechanical-waves:transport}
في \omterm{def:g10:signals-and-waves:wave}{موجة} الملعب يكون ``الوسط'' هو الجمهور: فالمنتقل هو
فعل الوقوف، بينما يحتفظ كل متفرج بمقعده. وعلى
البركة تتوسع الحلقة لكن فلّينة طافية لا تفعل إلا أن ترتفع وتنخفض في مكانها.
\end{example}

\begin{definition}[الأمواج المستعرضة والطولية]\label{def:g12:mechanical-waves:transverse}
تكون \omterm{def:g10:signals-and-waves:wave}{الموجة} \emph{مستعرضة}\index{موجة مستعرضة} إذا تحرك الوسط
\emph{عموديًّا} على الانتشار: فتنطلق حدبة على حبل
منفوض بينما لا يتحرك كل نقطة منه إلا صعودًا وهبوطًا. وتكون
\emph{طولية}\index{موجة طولية} إذا تحرك الوسط
\emph{على امتداد} الانتشار: فتنتقل منطقة من اللفّات المضغوطة على
زنبرك مدفوع، وكل لفّة تتأرجح في مكانها.
\end{definition}

\begin{omfigure}
\begin{tikzpicture}[font=\small]
  \node[left] at (-0.15,2.3) {مستعرضة};
  \draw[omDef, thick, smooth, samples=100, domain=0:8.4]
    plot (\x, {2.3 + 0.42*sin(130*\x)});
  \draw[omThm, Stealth-Stealth] (6.92,1.78) -- (6.92,2.82)
    node[right, yshift=-14pt] {الحركة};
  \draw[-Stealth, thick] (3.4,3.15) -- (5.0,3.15)
    node[midway, above] {الانتشار};
  \node[left] at (-0.15,0) {طولية};
  \foreach \i in {0,...,41}
    \draw[omDef, thick]
      ({0.2*\i + 0.26*sin(32*\i)},0.32) -- ({0.2*\i + 0.26*sin(32*\i)},-0.32);
  \node[below] at (3.4,-0.35) {تضاغط};
  \node[below] at (4.5,-0.35) {تخلخل};
  \draw[omThm, Stealth-Stealth] (0.75,0.55) -- (1.55,0.55)
    node[midway, above] {الحركة};
\end{tikzpicture}

{\small \omterm{def:g12:mechanical-waves:transverse}{المستعرضة} (الحبل): يتحرك الوسط عموديًّا على الانتشار.
\omterm{def:g12:mechanical-waves:transverse}{والطولية} (الزنبرك): تنتقل التضاغطات، وتتأرجح اللفّات في مكانها.}
\end{omfigure}

\begin{remark}[الصوت طولي]\label{rem:g12:mechanical-waves:sound}
\omterm{def:g10:signals-and-waves:sound}{الصوت} هو \omterm{def:g12:mechanical-waves:transverse}{الموجة الطولية} بامتياز: فالهواء المضغوط ---
أي \omterm{def:g10:pressure:pressure}{الضغط} الأعلى (\cref{ch:g10:pressure}) --- ينتقل من المنبع إلى
طبلة الأذن؛ وقد خُصّص له \cref{ch:g12:sound-acoustics}. أما الزلازل
فترسل النوعين معًا: فأمواج P تضغط الصخر، وأمواج S تهزّه جانبيًّا.
\end{remark}

\section{السرعة والتأخير}

\begin{definition}[سرعة الانتشار]\label{def:g12:mechanical-waves:speed}
\emph{سرعة انتشار}\index{سرعة الانتشار} موجةٍ هي
المسافة التي يقطعها الاضطراب مقسومة على الزمن المستغرَق:
$v = d/t$. وهي سرعة \emph{الشكل} المنتقل، لا سرعة
الوسط.
\end{definition}

\begin{omfigure}
\begin{tikzpicture}[font=\small]
  \draw[-Stealth] (-0.4,0) -- (8.9,0) node[below] {$x$ (\unit{m})};
  \foreach \x in {0,1,...,8}
    \draw (\x,0.05) -- (\x,-0.05) node[below] {$\x$};
  \draw[omDef, thick, smooth, samples=90, domain=-0.2:8.6]
    plot (\x, {1.1*exp(-3.5*(\x-1.0)^2)});
  \draw[omThm, thick, dashed, smooth, samples=90, domain=-0.2:8.6]
    plot (\x, {1.1*exp(-3.5*(\x-3.4)^2)});
  \node[omDef, above] at (1.0,1.15) {$t = 0$};
  \node[omThm, above] at (3.75,1.15) {$t = \qty{0.30}{s}$};
  \draw[omProp, Stealth-Stealth] (1.0,1.75) -- (3.4,1.75)
    node[midway, above] {$d = v\,\Delta t$};
\end{tikzpicture}

{\small نبضة على حبل مصوَّرة مرتين: انتقلت القمة \qty{2.4}{m}
في \qty{0.30}{s}، ومنه $v = \qty{8.0}{m/s}$ --- فالشكل ينتقل، و
الحبل يبقى.}
\end{omfigure}

\begin{proposition}[الوسط يحدّد السرعة]\label{prop:g12:mechanical-waves:medium}
في الاضطرابات الصغيرة، لا تتوقف \omterm{def:g12:mechanical-waves:speed}{سرعة الانتشار} إلا على
\emph{الوسط}، لا على شكل الاضطراب ولا على حجمه. فعلى
حبل مشدود، تزداد $v$ مع \omterm{def:g10:inertia:inventory}{الشدّ} وتقلّ مع الكتلة في
\omterm{def:g10:orders-of-magnitude:unit}{المتر}: فالمشدود الخفيف سريع، والمرتخي الثقيل بطيء. \admitted
\end{proposition}

\begin{remark}[أين تقيم الصيغة]\label{rem:g12:mechanical-waves:ropeformula}
تُشتقّ صيغة سرعة الحبل الدقيقة في مجلد السنة الأولى؛ وهذه
السنة تكفي الوقائع التجريبية أعلاه.
\end{remark}

\begin{example}[سرعات الصوت]\label{ex:g12:mechanical-waves:speeds}
ينتقل \omterm{def:g10:signals-and-waves:sound}{الصوت} بسرعة \qty{340}{m/s} في الهواء، و\qty{1500}{m/s} في الماء، و
\qty{5900}{m/s} في الفولاذ --- فكلما ازداد الوسط صلابة ازداد سرعة. ومن هنا
أذن أفلام الغرب الملتصقة بالسكة: فالفولاذ يعلن القطار أولًا.
\end{example}

\begin{definition}[التأخير]\label{def:g12:mechanical-waves:delay}
\omterm{def:g10:signals-and-waves:wave}{الموجة} التي سرعتها $v$ تبلغ نقطة على المسافة $d$ بعد
\emph{تأخير}\index{تأخير} قدره $\tau = d/v$: فنقطتان تفصل بينهما $d$
تكرّران الحركة نفسها، وتتأخر الأبعد بمقدار $\tau$ --- أي
القياس \omterm{rem:g10:signals-and-waves:honest}{بالصدى} في \cref{ch:g10:signals-and-waves}، مقروءًا معكوسًا.
\end{definition}

\begin{proposition}[التراكب بلا تفاعل]\label{prop:g12:mechanical-waves:superposition}
حيث تتراكب موجتان، يكون انتقال الوسط في كل
لحظة \emph{مجموع} الانتقالين اللذين تحدثهما كلٌّ منهما وحدها؛
وبعد التقاطع تواصل كل \omterm{def:g10:signals-and-waves:wave}{موجة} سيرها دون تغير --- وهو مشتقّ في
مجلد السنة الثانية من خطّية معادلات الحركة. \admitted
\end{proposition}

\begin{example}[نبضتان متقاطعتان]\label{ex:g12:mechanical-waves:crossing}
أرسِل حدبة صاعدة من كل طرف من حبل: فحيث تلتقيان، يرتفع
الحبل برهةً إلى مجموع الارتفاعين؛ ثم تخرج حدبتان سليمتان
--- كما تبلغ محادثتان متقاطعتان مستمعيهما.
\end{example}

\section{الأمواج الدورية: الدورية المزدوجة}

\begin{definition}[الموجة الدورية]\label{def:g12:mechanical-waves:periodic}
حين يكرّر المنبع حركته \omterm{def:g10:signals-and-waves:period}{بدور} $T$ \omterm{def:g10:signals-and-waves:frequency}{وتواتر}
$f = 1/T$ (\cref{ch:g10:signals-and-waves})، أنشأ في الوسط
\emph{موجة دورية}\index{موجة دورية}: فتكرّر كل نقطة حركتها
هي \omterm{def:g10:signals-and-waves:period}{بالدور} $T$، وتتأخر كلٌّ عن جيرانها \omterm{def:g12:mechanical-waves:delay}{بتأخير} الانتقال.
\end{definition}

\begin{definition}[طول الموجة]\label{def:g12:mechanical-waves:wavelength}
تُظهر صورة لحظية \omterm{def:g12:mechanical-waves:periodic}{لموجة دورية} نمطًا يتكرر في الفضاء؛ ودوره
المكاني --- من قمة إلى قمة --- هو
\emph{طول الموجة}\index{طول الموجة} $\lambda$، \omterm{def:g10:orders-of-magnitude:unit}{بالمتر}. \omterm{def:g12:mechanical-waves:periodic}{والموجة
الدورية} \emph{مزدوجة الدورية}\index{دورية مزدوجة}: \omterm{def:g10:signals-and-waves:period}{بدور} $T$
في الزمن عند كل نقطة، \omterm{def:g10:signals-and-waves:period}{ودور} $\lambda$ في الفضاء في كل لحظة.
\end{definition}

\begin{theorem}[علاقة الموجة]\label{thm:g12:mechanical-waves:relation}
\omterm{def:g12:mechanical-waves:periodic}{الموجة الدورية} التي سرعتها $v$ ودورها $T$ وتواترها $f$ طول
موجتها $\lambda = v\,T = v/f$.
\end{theorem}

\begin{proof}
خلال \omterm{def:g10:signals-and-waves:period}{دور} واحد ينزلق النمط إلى الأمام مسافة $v\,T$؛ لكن المنبع
--- ومن ثم كل نقطة --- يكون عندئذٍ قد عاد إلى حالته الابتدائية، فينطبق
النمط المنزلق على الأصلي: أي أن مسافة التكرار هي $v\,T$.
\end{proof}

\begin{omfigure}
\begin{tikzpicture}[font=\small]
  \draw[-Stealth] (-0.25,0) -- (4.35,0) node[below] {$x$};
  \draw[-Stealth] (0,-1.1) -- (0,1.55) node[left] {$y$};
  \draw[omDef, thick, smooth, samples=110, domain=0:4.1]
    plot (\x, {0.8*sin(160*\x)});
  \draw[omProp, Stealth-Stealth] (0.5625,1.05) -- (2.8125,1.05)
    node[midway, above] {$\lambda$};
  \node at (2.05,-1.55) {صورة لحظية $y(x)$، $t$ ثابت};
  \begin{scope}[xshift=5.6cm]
    \draw[-Stealth] (-0.25,0) -- (4.35,0) node[below] {$t$};
    \draw[-Stealth] (0,-1.1) -- (0,1.55) node[left] {$y$};
    \draw[omThm, thick, smooth, samples=110, domain=0:4.1]
      plot (\x, {0.8*sin(160*\x)});
    \draw[omProp, Stealth-Stealth] (0.5625,1.05) -- (2.8125,1.05)
      node[midway, above] {$T$};
    \node at (2.05,-1.55) {إشارة $y(t)$، $x$ ثابت};
  \end{scope}
\end{tikzpicture}

{\small الجيب نفسه مرتين: تتكرر الصورة اللحظية كل $\lambda$ في
الفضاء، وتتكرر \omterm{def:g10:signals-and-waves:signal}{إشارة} النقطة الواحدة كل $T$ في الزمن؛ و$\lambda = v\,T$.}
\end{omfigure}

\begin{remark}[التواتر يخصّ المنبع]\label{rem:g12:mechanical-waves:source}
عند الانتقال من وسط إلى آخر --- \omterm{def:g10:signals-and-waves:sound}{كالصوت} من الهواء إلى الماء ---
تحفظ \omterm{def:g10:signals-and-waves:wave}{الموجة} تواترها، الذي يمليه المنبع، بينما تتغير سرعتها؛
وبحسب $\lambda = v/f$ يتغير \omterm{def:g12:mechanical-waves:wavelength}{طول الموجة} معها.
\end{remark}

\begin{example}[نغمة لا، في الهواء وتحت الماء]\label{ex:g12:mechanical-waves:concertA}
نغمة لا عند \qty{440}{Hz} في \cref{ch:g10:signals-and-waves} طول
موجتها $340/440 = \qty{0.77}{m}$ في الهواء و$1500/440 =
\qty{3.4}{m}$ في الماء: فالنغمة نفسها، ممدودةً بالوسط الأسرع.
\end{example}

\begin{definition}[على توافق الطور وعلى تعاكسه]\label{def:g12:mechanical-waves:phase}
نقطتان تفصل بينهما عدد صحيح من أطوال \omterm{def:g10:signals-and-waves:wave}{الموجة}، $d = k\lambda$،
تتحركان تحركًا متطابقًا في كل لحظة: فهما
\emph{على توافق الطور}\index{توافق الطور}. أما النقطتان اللتان تفصل بينهما
$d = (k + \tfrac12)\lambda$ فتتحركان متعاكستين دائمًا --- قمة في مقابل
قاع: فهما على \emph{تعاكس الطور}\index{تعاكس الطور}.
\end{definition}

\begin{method}[قراءة الدورية المزدوجة]\label{met:g12:mechanical-waves:reading}
\begin{enumerate}
  \item على تسجيل $y(t)$ عند نقطة ثابتة، اقرأ $T$ من قمة
        إلى قمة؛ ومنه $f = 1/T$.
  \item وعلى صورة لحظية $y(x)$ عند لحظة ثابتة، اقرأ $\lambda$ من قمة
        إلى قمة.
  \item استنتج $v = \lambda/T = \lambda f$؛ وتحقّق من ذلك بتوقيت
        مباشر $v = d/\tau$ متى توفّر.
  \item ولا تخلط بين المخططين أبدًا: فمحور أحدهما يحمل الثواني، ومحور
        الآخر يحمل الأمتار.
\end{enumerate}
\end{method}

\section{جبهات الموجة والتوهين}

\begin{definition}[جبهة الموجة]\label{def:g12:mechanical-waves:wavefront}
\emph{جبهة الموجة}\index{جبهة الموجة} خط (أو سطح) من النقاط
تبلغها \omterm{def:g10:signals-and-waves:wave}{الموجة} في اللحظة نفسها --- كخط قمة على
البركة. ومن منبع نقطي تكون جبهات \omterm{def:g10:signals-and-waves:wave}{الموجة} دوائر متوسعة؛ وبعيدًا
عن أي منبع تكون شبه مستقيمة. وتنتشر \omterm{def:g10:signals-and-waves:wave}{الموجة}
عموديًّا على جبهاتها، والقمم المتتالية تفصل بينها $\lambda$.
\end{definition}

\begin{definition}[التوهين]\label{def:g12:mechanical-waves:attenuation}
تناقص \omterm{def:g10:signals-and-waves:amplitude}{سعة} \omterm{def:g10:signals-and-waves:wave}{موجة} أثناء انتشارها هو
\emph{توهينها}\index{توهين}. وسببان يتجمعان: \omterm{def:g10:energy-conservation:energy}{فطاقة}
جبهة \omterm{def:g10:signals-and-waves:wave}{موجة} دائرية تُتقاسم على قمة تزداد طولًا، و
الوسط يمتصّ قليلًا في العبور. أما الجبهات المستقيمة فتفلت من السبب
الأول --- وبذلك يعبر موج طويل لعاصفة بعيدة محيطًا.
\end{definition}

\begin{omfigure}
\begin{tikzpicture}[font=\small]
  \fill (0,0) circle (2pt);
  \foreach \r/\c in {0.5/100, 1.0/82, 1.5/64, 2.0/48, 2.5/34}
    \draw[thick, omDef!\c!white] (0,0) circle (\r);
  \draw[omProp, Stealth-Stealth] (315:1.5) -- (315:2.0)
    node[midway, below left, inner sep=1pt] {$\lambda$};
  \draw[-Stealth, omThm, thick] (25:2.55) -- (25:3.25)
    node[right] {$v$};
  \foreach \x/\c in {4.9/100, 5.5/90, 6.1/80, 6.7/70}
    \draw[thick, omDef!\c!white] (\x,-1.5) -- (\x,1.5);
  \draw[-Stealth, omThm, thick] (6.9,0) -- (7.6,0) node[right] {$v$};
\end{tikzpicture}

{\small جبهات \omterm{def:g10:signals-and-waves:wave}{موجة} دائرية، يفصل بينها طول \omterm{def:g10:signals-and-waves:wave}{موجة} واحد، تخبو كلما انتشرت
\omterm{def:g10:energy-conservation:energy}{الطاقة} على قمم أطول؛ أما الجبهات المستقيمة (يسارًا) فتخبو ببطء.}
\end{omfigure}

\section{تمارين}

\begin{exercise}[$\star$]\label{exo:g12:mechanical-waves:1}
تجري \omterm{def:g10:signals-and-waves:wave}{موجة} ملعب على حلقة المدرجات البالغة \qty{320}{m} في \qty{40}{s}.
(أ) احسب سرعتها. (ب) وهل يدور أي متفرج حول الحلقة؟
فما الذي يدور؟ (ج) يقف كل متفرج بعد \qty{1.5}{s} من جاره
الذي يبعد \qty{12}{m}: تحقّق من الاتساق.
\end{exercise}

\begin{exercise}[$\star$]\label{exo:g12:mechanical-waves:2}
\omterm{def:g12:mechanical-waves:transverse}{مستعرضة} أم \omterm{def:g12:mechanical-waves:transverse}{طولية}؟ (أ) نبضة على حبل يُنفض صعودًا وهبوطًا؛
(ب) تضاغط يُرسل على زنبرك؛ (ج) \omterm{def:g10:signals-and-waves:sound}{الصوت} في الهواء؛ (د) تموّج
بركة (السطح يرتفع وينخفض عموديًّا، والحلقة تنتقل أفقيًّا)؛ (ه)
أمواج P الزلزالية (تضغط في جهة الانتقال) وأمواج S (تهزّ جانبيًّا).
\end{exercise}

\begin{exercise}[$\star$]\label{exo:g12:mechanical-waves:3}
يبلغك الرعد بعد \qty{4.0}{s} من الومضة. (أ) فكم تبعد
الصاعقة ($v = \qty{340}{m/s}$)؟ (ب) ويسمع غطّاس \omterm{def:g12:mechanical-waves:delay}{التأخير} نفسه بين
ومضة ودويّ تحت الماء ($v = \qty{1500}{m/s}$): فكم يبعد ذلك
المنبع؟ (ج) وما الذي يجب أن تعرفه قبل تحويل \omterm{def:g12:mechanical-waves:delay}{التأخير} إلى مسافة؟
\end{exercise}

\begin{exercise}[$\star$]\label{exo:g12:mechanical-waves:4}
احسب طول \omterm{def:g10:signals-and-waves:wave}{موجة} \omterm{def:g10:signals-and-waves:sound}{صوت} تواتره \qty{440}{Hz} في الهواء وفي الماء.
وحين يعبر \omterm{def:g10:signals-and-waves:sound}{الصوت} من الهواء إلى الماء، أيّ المقادير $f$ و$v$
و$\lambda$ يتغير، وأيّها يبقى ثابتًا؟
\end{exercise}

\begin{exercise}[$\star$]\label{exo:g12:mechanical-waves:5}
على بركة، تفصل بين قمم التموّج المتجاورة \qty{12}{cm} وتتقدم الحلقات
بسرعة \qty{0.30}{m/s}. جِد \omterm{def:g12:mechanical-waves:wavelength}{طول الموجة} \omterm{def:g10:signals-and-waves:period}{والدور} و
\omterm{def:g10:signals-and-waves:frequency}{التواتر}.
\end{exercise}

\begin{exercise}[$\star\star$]\label{exo:g12:mechanical-waves:6}
تُصوَّر نبضة على حبل طويل مرتين، يفصل بينهما \qty{0.30}{s}: فتقع
قمتها عند $x = \qty{1.0}{m}$، ثم عند \qty{3.4}{m}. (أ) فما السرعة؟
(ب) ومتى تبلغ النبضة الجدار عند $x = \qty{7.0}{m}$؟ (ج) وشريط معقود
عند $x = \qty{5.0}{m}$: فما حركته عند مرور النبضة؟
\end{exercise}

\begin{exercise}[$\star\star$]\label{exo:g12:mechanical-waves:7}
نبضتان صاعدتان، ارتفاعاهما \qty{3.0}{cm} و\qty{2.0}{cm}، تنتقلان
إحداهما نحو الأخرى بسرعة \qty{2.0}{m/s} لكلٍّ، والقمتان تفصل بينهما \qty{4.0}{m}.
(أ) فمتى تتطابق القمتان؟ (ب) وما أقصى ارتفاع للحبل عندئذٍ؟
(ج) والسؤالان نفساهما لو كانت النبضة الصغيرة نازلة. (د) وما حال الحبل بعد ذلك مباشرة؟
\end{exercise}

\begin{exercise}[$\star\star$]\label{exo:g12:mechanical-waves:8}
(أ) حبلا غسيل يحملان \omterm{def:g10:inertia:inventory}{الشدّ} نفسه؛ وتُرسَل نبضة على
كلٍّ منهما. فأيّهما أسرع: الرفيع أم الغليظ؟ ولماذا؟ (ب) ويشدّ
عازف قيثارة وترًا: فماذا يحدث لسرعة \omterm{def:g10:signals-and-waves:wave}{الموجة}؟ (ج) وهل يجعل النفض
بضعف \omterm{def:g10:inertia:force}{القوة} (أي نبضة أعلى) النبضةَ تنتقل أسرع؟
\end{exercise}

\begin{exercise}[$\star\star$]\label{exo:g12:mechanical-waves:9}
يضرب عامل سكة على بعد \qty{1.8}{km}؛ وأنت تصغي بأذن على
الفولاذ. احسب زمني الوصول عبر السكة (\qty{5900}{m/s})
وعبر الهواء (\qty{340}{m/s})، والفاصل بين الطرقتين.
\end{exercise}

\begin{exercise}[$\star\star$]\label{exo:g12:mechanical-waves:10}
يحمل حبل \omterm{def:g12:mechanical-waves:periodic}{موجة دورية} طول موجتها \qty{24}{cm}. فهل النقطتان \omterm{def:g12:mechanical-waves:phase}{على توافق الطور}
أو على تعاكسه أو لا هذا ولا ذاك، عند تباعد قدره (أ) \qty{48}{cm}؛
(ب) \qty{12}{cm}؛ (ج) \qty{36}{cm}؛ (د) \qty{30}{cm}؛ (ه) \qty{6.0}{cm}؟
\end{exercise}

\begin{exercise}[$\star\star$]\label{exo:g12:mechanical-waves:11}
يرسل زلزال أمواج P بسرعة \qty{6.0}{km/s} وأمواج S بسرعة
\qty{3.5}{km/s}، ودورهما \qty{2.0}{s}. احسب طولي \omterm{def:g10:signals-and-waves:wave}{الموجة}،
وقارنهما ببناية، واشرح لماذا يرتفع حيّ بأكمله وينخفض
معًا تقريبًا.
\end{exercise}

\begin{exercise}[$\star\star\star$]\label{exo:g12:mechanical-waves:12}
يُظهر تسجيل ارتفاع الماء عند عوّامة قممًا كل
\qty{0.50}{s}؛ وتُظهر صورة جوية قممًا تفصل بينها \qty{0.75}{m}.
(أ) جِد $T$ و$f$. (ب) جِد $\lambda$. (ج) استنتج السرعة. (د) وتركب
العوّامة قمةً عند $t = 0$: فمتى تركب القمة التالية؟ (ه) وهل نقطة تبعد \qty{1.875}{m}
\omterm{def:g12:mechanical-waves:phase}{على توافق الطور} مع العوّامة، أم على تعاكسه، أم لا هذا ولا ذاك؟
\end{exercise}

\begin{exercise}[$\star\star\star$]\label{exo:g12:mechanical-waves:13}
تكمل عوّامة راسية $10$ تذبذبة كاملة في \qty{80}{s} على
موج طويل منتظم تفصل بين قممه \qty{120}{m}. (أ) جِد \omterm{def:g10:signals-and-waves:period}{الدور} و
\omterm{def:g10:signals-and-waves:frequency}{التواتر}. (ب) جِد سرعة الموج، \omterm{def:g10:orders-of-magnitude:unit}{بوحدة} \unit{m/s} وبوحدة
\unit{km/h}. (ج) وكم تستغرق قمة عند العوّامة الآن لتبلغ
الشاطئ الذي يبعد \qty{900}{m}؟ (د) وهل تنجرف العوّامة إلى الشاطئ معها؟
\end{exercise}

\begin{exercise}[$\star\star\star$]\label{exo:g12:mechanical-waves:14}
نبضتان متطابقتان، إحداهما صاعدة والأخرى نازلة، تنتقلان إحداهما نحو
الأخرى بسرعة \qty{3.0}{m/s} لكلٍّ، ومركزاهما يفصل بينهما \qty{6.0}{m}. (أ) فمتى
يتطابق مركزاهما، وكيف يبدو الحبل عندئذٍ؟ (ب) وهل يكون
الحبل ساكنًا عندئذٍ؟ وأين ذهبت \omterm{def:g10:energy-conservation:energy}{الطاقة}؟ (ج) وماذا يحدث
تاليًا؟ (د) ولماذا لا يعني ``مسطّح للحظة'' أن ``لا شيء هناك''؟
\end{exercise}

\begin{exercise}[$\star\star\star$]\label{exo:g12:mechanical-waves:15}
ينتشر تموّج حصاة على بركة ساكنة. (أ) احسب طول
القمة عند $r = \qty{0.50}{m}$ وعند $r = \qty{8.0}{m}$؛ فما عامل النمو؟
(ب) وتُحفظ \omterm{def:g10:energy-conservation:energy}{طاقة} القمة تقريبًا: فماذا يحدث \omterm{def:g10:energy-conservation:energy}{للطاقة}
في \omterm{def:g10:orders-of-magnitude:unit}{المتر}، وللسعة؟ (ج) ولماذا يتوهّن موج طويل ذو جبهة مستقيمة
أقلّ بكثير؟ (د) ولماذا لا تتوهّن \omterm{def:g10:signals-and-waves:wave}{موجة} الملعب البتة؟
\end{exercise}

\section{مسألة: قراءة مخطط زلزالي}

\begin{problem}[{مسألة نهاية الأسبوع --- قراءة مخطط زلزالي: موجتان
تنطلقان من صدع مكسور، وثلاث دوائر تحدّد مركز الزلزال، وموجة ثالثة
بطيئة تعطي ساحلًا إنذاره}]
\label{pb:g12:mechanical-waves:1}
مناوبة ليلية في مرصد للزلازل. عند الساعة 09:14:32 تقفز إبرة
المحطة A: تموّجات سريعة حادّة، ثم تأرجحات أبطأ أوسع. لقد
كسر زلزال صخرًا في مكان ما؛ وتنطلق موجتان عبر
القشرة: أمواج P، وهي تضاغطات بسرعة $v_P = \qty{6.0}{km/s}$، وأمواج S،
وهي هزّ جانبي بسرعة $v_S = \qty{3.5}{km/s}$.

\medskip
\noindent\textbf{الجزء I --- موجتان وساعة واحدة.}
\begin{enumerate}
  \item انطلاقًا من الأوصاف أعلاه، صنّف أمواج P وS
        بين \omterm{def:g12:mechanical-waves:transverse}{مستعرضة} \omterm{def:g12:mechanical-waves:transverse}{وطولية}.
  \item وأيّ \omterm{def:g10:signals-and-waves:wave}{موجة} تصل أولًا إلى كل محطة؟ احسب زمني
        انتقالهما على $d = \qty{120}{km}$.
  \item بيّن أن الفارق بين وصول P وS هو
        $\Delta t = d\,(1/v_S - 1/v_P)$ واحسبه من أجل
        $d = \qty{120}{km}$.
  \item استنتج قاعدة قياس المسافة عند عالم الزلازل
        $d = \dfrac{v_P\,v_S}{v_P - v_S}\,\Delta t$ وبيّن أن
        \omterm{def:g10:orders-of-magnitude:scinot}{المعامل} نحو \qty{8.4}{km} لكل \omterm{def:g10:orders-of-magnitude:unit}{ثانية} من الفارق.
  \item ولماذا ينمو الفارق مع المسافة --- ولماذا تكون تلك هدية
        لمحطة لم تشهد بداية الزلزال قط؟
\end{enumerate}

\medskip
\noindent\textbf{الجزء II --- دوائر على الخريطة.}
\begin{enumerate}[resume]
  \item تقيس المحطة A فارقًا قدره $\Delta t_A = \qty{25}{s}$. فكم يبعد
        مركز الزلزال عن A؟
  \item اشرح لماذا تضع هذه القراءة الواحدة مركز الزلزال على دائرة،
        لا على نقطة.
  \item وتقيس المحطتان B وC فارقين قدرهما $\Delta t_B = \qty{15}{s}$ و
        $\Delta t_C = \qty{32}{s}$: فما مسافتاهما إلى مركز الزلزال؟
  \item اشرح كيف تحدّد الدوائر الثلاث معًا مركز الزلزال.
  \item بلغت \omterm{def:g10:signals-and-waves:wave}{موجة} P المحطة A عند 09:14:32. احسب زمن انتقالها من
        مركز الزلزال، وزمن نشوء الزلزال.
\end{enumerate}

\medskip
\noindent\textbf{الجزء III --- \omterm{def:g10:signals-and-waves:wave}{الموجة} التي تعبر المحيط.}
تلتقي الدوائر الثلاث في عرض البحر: فقد انكسر الصدع تحت قاع البحر
وحرّك عمود الماء. والأمواج الطويلة على الماء العميق تنتقل
بسرعة يحدّدها العمق --- نحو \qty{200}{m/s} في المحيط المفتوح،
ونحو \qty{10}{m/s} قرب شاطئ (نقبله هنا؛ وتُشتقّ الصيغة
في مجلد السنة الأولى). ويقع ميناء على بعد $D = \qty{600}{km}$ من
مركز الزلزال.
\begin{enumerate}[resume]
  \item حوّل \qty{200}{m/s} إلى \unit{km/h} وسمِّ مركبةً
        سرعة إبحارها مقاربة.
  \item وكم يستغرق التسونامي ليبلغ الميناء؟
  \item ويحسّ مقياس الزلازل في الميناء \omterm{def:g10:signals-and-waves:wave}{موجة} P أولًا. فبعد كم
        من الوقت؟ وإذا انطلق الإنذار في تلك اللحظة، فكم زمنًا من
        الإنذار يناله الساحل قبل التسونامي؟
  \item \omterm{def:g10:signals-and-waves:period}{ودور} التسونامي في البحر نحو \qty{15}{min}. احسب
        \omterm{def:g10:light-spectra:wavelength}{طول موجته} هناك، واشرح لماذا لا تلاحظه سفينة.
  \item وباقترابه من الشاطئ تهبط السرعة إلى \qty{10}{m/s} بينما
        يبقى \omterm{def:g10:signals-and-waves:period}{الدور} ثابتًا. احسب \omterm{def:g12:mechanical-waves:wavelength}{طول الموجة} الجديد. فماذا يجب أن
        يحدث \omterm{def:g10:signals-and-waves:wave}{للموجة} وهي تتكدّس؟
\end{enumerate}

\medskip
\noindent\textbf{الجزء IV --- \omterm{def:g10:energy-conservation:energy}{الطاقة}، والتقرير.}
يحرّر زلزال قوي من رتبة \qty{1e15}{J}، تحملها
الأمواج.
\begin{enumerate}[resume]
  \item احسب طول جبهة \omterm{def:g10:signals-and-waves:wave}{موجة} دائرية على بعد \qty{10}{km} من
        مركز الزلزال، ثم على بعد \qty{210}{km}؛ فبأي عامل نمت؟
  \item استنتج لماذا يضعف الهزّ مع المسافة حتى في قشرة
        بلا ضياع --- وسمِّ السبب الثاني الذي تضيفه قشرة \omterm{def:g11:lenses-and-eye:images}{حقيقية}.
  \item السائل لا يقاوم القصّ. وتسجّل محطات الجهة البعيدة \omterm{def:g10:signals-and-waves:wave}{موجة}
        P ولا تسجّل \omterm{def:g10:signals-and-waves:wave}{موجة} S: فماذا يكشف ذلك عن \omterm{def:g11:fundamental-interactions:ladder}{نواة} الأرض؟
  \item القطار السريع في فصل \omterm{def:g10:energy-conservation:energy}{الطاقة} للسنة \omterm{def:g10:orders-of-magnitude:prefixes}{السابقة}
        (\cref{ch:g11:mechanical-energy}) يحمل نحو \qty{1.5e9}{J}
        من \omterm{def:g11:mechanical-energy:kinetic}{الطاقة الحركية}: فكم قطارًا كهذا يساوي \qty{1e15}{J}؟
  \item التقرير، في جملتين: أين ومتى وقع الزلزال،
        وكم إنذارًا نال الميناء؟ أعطِ الأعداد الثلاثة.
\end{enumerate}
\end{problem}
